\chapter*{Abstract}
\addcontentsline{toc}{chapter}{Abstract}

% Written while the project goes on: it describes the goal and the state of the work, and will be
% rewritten with the results at the end (about 200-300 words).
Production scheduling assigns activities to resources over time while respecting precedences,
capacities, calendars and deadlines. The same problem appears well beyond the factory floor, for
instance in public administration offices that process files through a sequence of steps.
Exact solvers find optimal plans but their computation time grows quickly with the size of the
instance, while classic dispatching rules are fast but myopic.

This document describes \emph{FlowSync}, an AI-assisted scheduling platform that a team of three
students is building around the idea that the AI proposes the plan and the planner decides. In
the intended design, a request, possibly written in natural language, is turned into a
structured problem specification; a first model (M1) estimates uncertain durations from the
history of the executions, and a second model (M2) builds the schedule step by step inside a
simulator, so that every plan is feasible by construction. M2 is to be trained by imitation of a
constraint programming solver (CP-SAT) used only offline. For each request the planner receives
three alternative plans, compared with the ATC and EDD dispatching rules.

The document is written while the project goes on, and each chapter states how far the work has
come. At the time of writing, the architecture and the contracts between the parts are defined;
the data pipeline of the AI (instance generator, dispatcher, dispatching rules, CP-SAT teacher
and training dataset) is implemented and evaluated; the backend provides the database,
authentication, master data and orders; and the web interface is implemented against a mock of
the backend. The models, the AI service and the planning functions of the backend are the next
steps.

\vspace{0.5cm}
\noindent\textbf{Keywords:} production scheduling, imitation learning, dispatching rules,
constraint programming, decision support systems.
